% ============================================================
%  SCHOLA DOMESTICA — Level 1, Week 4, Lesson 19
%  Topic: Numbers 0–10 — Full Review & Counting Backward from 10
%  10–12 problems | No speed drill
%  Compile with XeLaTeX
% ============================================================
\documentclass[12pt,letterpaper]{article}
\usepackage{sd_style}

\renewcommand{\lessonnum}{19}
\renewcommand{\lessonweek}{4}

\begin{document}

\lessonheader{4}{19}{Numbers 0–10: Full Review \& Counting Backward}

\begin{instructbox}
\noindent\textbf{\color{sd-blue}Instruction} \textit{(read aloud):}\\[4pt]
Today we review all numbers 0 through 10 and practice counting
backward from 10 all the way to 0.
Counting backward from 10 is like a grand rocket launch!
\end{instructbox}

\vspace{6pt}
\begin{center}
\begin{tikzpicture}
  \draw[sd-blue, very thick, <-] (-0.3,0) -- (10.8,0);
  \foreach \n in {0,...,10}{
    \draw[sd-blue, thick] (\n, -0.15) -- (\n, 0.15);
    \node[below] at (\n, -0.25) {\textbf{\n}};
  }
  \draw[sd-gold, very thick, ->] (10.2, 0.5) -- (9.0, 0.5);
  \node[above] at (9.6, 0.55) {\color{sd-gold}\small\textit{count this way}};
\end{tikzpicture}
\end{center}

\goldrule
\noindent\textbf{\color{sd-blue}Practice}
\vspace{6pt}

\prob Count backward. Fill in the blanks.\\[6pt]
\hspace{1cm}{\Large 10, \ansblank, 8, \ansblank, 6, \ansblank, 4, \ansblank, 2, \ansblank, 0}

\vspace{10pt}
\prob What is \textbf{one less} than each number?\\[6pt]
\begin{tabular}{ccccc}
one less than 10: & one less than 7: & one less than 4: & one less than 9: & one less than 1:\\[4pt]
\ansblank & \ansblank & \ansblank & \ansblank & \ansblank
\end{tabular}

\vspace{10pt}
\prob Count and write.\\[8pt]
\begin{tabular}{ccccc}
\onedot\onedot\onedot\onedot\onedot\onedot\onedot\onedot\onedot &
\onedot\onedot\onedot\onedot\onedot\onedot &
\onedot\onedot\onedot\onedot\onedot\onedot\onedot\onedot\onedot\onedot &
\onedot\onedot\onedot\onedot\onedot\onedot\onedot &
\onedot\onedot\onedot\onedot\onedot\onedot\onedot\onedot\\[4pt]
\ansblank & \ansblank & \ansblank & \ansblank & \ansblank
\end{tabular}

\vspace{10pt}
\prob Circle every number that is \textbf{less than 5}:\\[6pt]
\hspace{2cm}{\Large 7 \quad 2 \quad 9 \quad 0 \quad 4 \quad 6 \quad 3 \quad 10 \quad 1}

\vspace{10pt}
\prob Write the missing number in each sequence:\\[6pt]
\hspace{1cm}0, 1, \ansblank, 3, 4 \hspace{2cm}
6, 7, \ansblank, 9, 10 \hspace{2cm}
10, 9, \ansblank, 7, 6

\vspace{10pt}
\prob Order from greatest to least:\\[6pt]
\hspace{1cm}{\Large 3 \quad 10 \quad 1 \quad 8 \quad 5}\quad$\longrightarrow$\quad
\ansblank\quad\ansblank\quad\ansblank\quad\ansblank\quad\ansblank

\vspace{10pt}
\prob \textit{(Teacher reads)} Ten candles were lit on a birthday cake.
The birthday child blew out one candle at a time.
Count backward to show how many were left after each blow:\\[6pt]
\hspace{1cm}10, \ansblank, \ansblank, \ansblank, \ansblank, \ansblank, \ansblank, \ansblank, \ansblank, \ansblank, \ansblank

\vspace{10pt}
\prob What number is exactly halfway between 0 and 10? \quad\ansblank

\vspace{10pt}
\prob Write all numbers from 0 to 10 in your neatest hand:\\[8pt]
\noindent\fbox{\rule{0pt}{1.8cm}\hspace{12.5cm}}

\goldrule

\begin{checkbox}
\noindent\textbf{\color{sd-green}Knowledge Check}\\[4pt]
\setcounter{probnum}{0}
\prob Count backward from 10 to 0 aloud. \quad $\square$~Done\\[8pt]
\prob What is one less than 6? \quad\ansblank\\[8pt]
\prob What is one less than 10? \quad\ansblank
\end{checkbox}

\vfill
\noindent{\color{sd-blue}$\bigstar$ \textbf{Enrichment — Cuneiform Numbers:}}\\[4pt]
The ancient Sumerians of Mesopotamia wrote numbers using a wedge-shaped stylus pressed into clay.
Their symbols for 1 through 10 looked like this:\\[4pt]
\begin{center}
\begin{tabular}{cccccccccc}
$\triangledown$ & $\triangledown\triangledown$ & $\triangledown\triangledown\triangledown$ &
$\triangledown\triangledown\triangledown\triangledown$ &
$\triangledown\triangledown\triangledown\triangledown\triangledown$ &
$\triangledown\triangledown\triangledown\triangledown\triangledown\triangledown$ &
$\triangledown\triangledown\triangledown\triangledown\triangledown\triangledown\triangledown$ &
$\triangledown\triangledown\triangledown\triangledown\triangledown\triangledown\triangledown\triangledown$ &
$\triangledown\triangledown\triangledown\triangledown\triangledown\triangledown\triangledown\triangledown\triangledown$ &
$\langle$\\[2pt]
1&2&3&4&5&6&7&8&9&10
\end{tabular}
\end{center}
\textit{(The wedge $\triangledown$ = 1; the angle $\langle$ = 10.)}

\vspace{4pt}\hrule\vspace{3pt}
{\footnotesize\color{gray}\textit{Teacher note: Counting backward from 10 is harder than from 5 for most children.
  Practice orally each day this week. The birthday candle story problem
  makes the concept memorable.}}

\end{document}
